%
\section{Node Orderings for Curved Geometry}
\label{sec:spatialdomains-orderings}

Section \ref{sec:spatialdomains-fundamentals} described how curvature
information reaches a geometry object: as a list of physical-space point
positions sampled at a known distribution of points on the reference element,
from which the modal coefficients of the mapping $\chi(\vec{\xi})$ are
inferred.  That list lives in a \code{Curve} object, held by
\code{Geometry::m\_curve}, and consists of a \code{PointsType} naming the
distribution together with a flat \code{std::vector} of \code{PointGeom}
pointers.

Being a flat vector, the \code{Curve} carries no information about which node
sits where on the reference element -- that is implied entirely by the
\code{PointsType} and by the shape of the geometry.  Getting this ordering
wrong does not produce an error; it produces a geometry that is subtly and
silently the wrong shape.  This section therefore documents, shape by shape,
what the ordering actually is.

The two places that have to agree on it are:

\begin{itemize}
\item \code{v\_FillGeom()}, which reads a \code{Curve} and interpolates it onto
  the quadrature points of \code{m\_xmap} to recover the mapping
  coefficients; and
\item \code{v\_MakeOrder()}, which does the reverse -- it evaluates the
  existing mapping and writes out a \code{Curve} at a new (usually higher)
  polynomial order.
\end{itemize}

\noindent Both are also the contract with any external mesh source, since a
\code{Curve} assembled by \code{MeshGraphIO} or by \nek{}'s mesh generator
\code{NekMesh} must be laid out the same way.

\subsection{Two Orderings: Modal and Nodal}
\label{sec:spatialdomains-orderings-two}

There are \emph{two} distinct node orderings in use within \nek{}, and the
difference between them is a frequent source of confusion.  It is deliberate,
and is a consequence of the order in which the code was written.

\begin{description}
\item[The modal ordering] is the one described by the \code{edgeVerts} and
  \code{faceVerts} tables inside the shape-specific geometry classes
  (\code{TetGeom}, \code{PrismGeom}, and so on).  It is the ordering that
  falls out naturally from the tensor-product construction of the expansion
  bases, and is what the vertex, edge, face and interior \emph{modes} of an
  expansion are numbered against.  It was revised at one point precisely so as
  to line up with the tensor-product structure.

\item[The nodal ordering] is the ordering emitted by the \code{eNodal*} point
  distributions in LibUtilities (for example
  \code{eNodalTetEvenlySpaced}), and consumed by the nodal expansions
  \code{StdNodalTriExp}, \code{StdNodalTetExp} and \code{StdNodalPrismExp}.
  It predates the modal ordering above.  \code{NekMesh} was written against
  it, and so the high-order node ordering in \nek{}'s mesh formats follows the
  nodal convention.
\end{description}

The practical consequence is that for the simplicial and prismatic shapes,
whose curves use a nodal distribution, the layout of a \code{Curve} is
\emph{not} simply the shape's own \code{edgeVerts} and \code{faceVerts}
tables read in order.  Assuming that it is will give you an element that is
correct at the vertices and wrong in the interior.  The known discrepancies
are called out per shape below; do not assume that others are absent.

A third layer is worth distinguishing from both of the above.  Sections
\ref{sect:node_ordering} and \ref{sect:HO_node_ordering} document the
ordering \code{NekMesh} uses internally and in \nek{}'s mesh formats, which
is a plain concatenation of vertices, then mid-edge nodes, then mid-face
nodes, then mid-volume nodes, for \emph{every} shape.  That is the nodal
convention, and for the simplicial and prismatic shapes it is also the
\code{Curve} layout.  For the tensor-product shapes it is not: a
quadrilateral or hexahedral \code{Curve} is a tensor grid, and the conversion
happens on the way in, in \code{GetCurvedNodesQuad()} and its relatives in
\code{NekMesh/MeshElements/}.  So a concatenated node list from a mesh file
and the \code{Curve} that ends up on a \code{QuadGeom} are two different
orderings of the same nodes.

\subsection{The FillGeom / MakeOrder Contract}
\label{sec:spatialdomains-orderings-contract}

Curvature is shared: an edge node belongs to the \code{SegGeom} for that edge,
and every element around that edge must see the same \code{PointGeom} object,
or the mesh will be discontinuous.  \code{v\_MakeOrder()} therefore does not
evaluate boundary nodes on its own mapping.  It \emph{copies} them from the
curves of its attached edges and faces, and evaluates only the nodes it
genuinely owns -- the interior ones.  This has two implications.

First, order of operations matters, and \code{MakeOrder} must be applied
bottom-up: all edges, then all faces, then all elements.  Each
\code{v\_MakeOrder()} asserts that the curves it needs are already present.

Second, and less obviously, \textbf{\code{MakeOrder} must never call
\code{FillGeom}}.  It evaluates the mapping that is already in place, which
means an unfilled geometry silently yields nodes at the origin; the temptation
to close that hole by filling on entry should be resisted.  If an attached edge
or face has already been modified, filling at that point would pull the
modification into the parent and subtly move the geometry part-way through the
elevation.  The protocol is instead two separate passes over the whole mesh:
\textbf{everyone calls \code{FillGeom}, and then everyone calls
\code{MakeOrder}}.  Callers are responsible for the first pass.

\subsection{Per-Shape Layouts}
\label{sec:spatialdomains-orderings-shapes}

Throughout, $P$ denotes the number of points along an edge, so that
$P = \text{order} + 1$, and the reference coordinates are
$\vec{\xi} = (\xi_1, \xi_2, \xi_3)$.

\subsubsection{SegGeom}

A one-dimensional distribution of $P$ points, running from vertex $0$ to
vertex $1$:
%
$$
[\, \vec{v}_0,\ \text{interior} \ldots,\ \vec{v}_1 \,].
$$
%
The end points are the vertices themselves, so the distribution must be one
that includes its end points (a GLL or evenly-spaced distribution, not a Gauss
one).

\subsubsection{QuadGeom}

A full tensor grid of $P^2$ points.  Grid point $(i,j)$ sits at
$\vec{\xi} = (\code{px}[j], \code{px}[i])$ and is stored at index
$i P + j$; that is, $\xi_1$ varies fastest.

The boundary of the grid is copied from the four edge curves, which means the
vertices come across for free as the ends of those edges.  Only the
$(P-2)^2$ interior nodes are evaluated.  The four element-local edges traverse
the grid as follows, with $k$ running along the local edge direction:

\begin{center}
\begin{tabular}{llll}
\hline
Edge & Constraint & $i$ & $j$ \\
\hline
0 & $\xi_2 = -1$, $\xi_1$ increasing & $0$     & $k$ \\
1 & $\xi_1 = +1$, $\xi_2$ increasing & $k$     & $P-1$ \\
2 & $\xi_2 = +1$, $\xi_1$ decreasing & $P-1$   & $P-1-k$ \\
3 & $\xi_1 = -1$, $\xi_2$ decreasing & $P-1-k$ & $0$ \\
\hline
\end{tabular}
\end{center}

\paragraph{Trap.}  \code{QuadGeom}'s constructor \emph{inverts}
\code{m\_eorient[2]} and \code{m\_eorient[3]} after computing them, so that
those stored values describe the edge relative to the anticlockwise circuit of
the boundary rather than relative to the element's local edge direction.  Code
that needs the latter -- as \code{v\_MakeOrder()} does, to decide whether to
read an edge curve forwards or backwards -- must flip them back.

\subsubsection{TriGeom}

A nodal distribution of $P(P+1)/2$ points, laid out as three vertices, then
the $P-2$ interior nodes of each of the three edges in \code{edgeVerts} order,
then the face interior:
%
$$
[\, \underbrace{\vec{v}_0, \vec{v}_1, \vec{v}_2}_{3},\
     \underbrace{e_0, e_1, e_2}_{3(P-2)},\
     \underbrace{\text{interior}}_{(P-2)(P-3)/2} \,].
$$

\paragraph{Trap.}  As with \code{QuadGeom}, the constructor inverts
\code{m\_eorient[2]}, and \code{v\_MakeOrder()} flips it back.

\subsubsection{HexGeom}

A full tensor grid of $P^3$ points.  Grid point $(d_0, d_1, d_2)$ sits at
$\vec{\xi} = (\code{px}[d_0], \code{px}[d_1], \code{px}[d_2])$ and is stored
at index
%
$$
d_0 + P\,(d_1 + P\,d_2),
$$
%
matching the ordering that \code{v\_FillGeom()} hands to \code{Interp3D}.

Because a quadrilateral's curve is itself a full tensor grid including its own
edge and vertex nodes, copying the six face curves brings across every
boundary node of the hexahedron, and only the $(P-2)^3$ interior nodes need to
be evaluated.  Each face's own $(a,b)$ indexing is related to the element's
view of that face by \code{StdExpansion::ReOrientTracePhysMap}, which for a
face orientation returns \code{idmap} such that \code{idmap[dest]} is the
source index; the resulting $(a,b)$ then maps onto element directions
following \code{v\_GetDir()} and the \code{faceVerts} table:

\begin{center}
\begin{tabular}{llll}
\hline
Face & Constraint & Spans & $(d_0, d_1, d_2)$ \\
\hline
0 & $\xi_3 = -1$ & $(\xi_1, \xi_2)$ & $(a,\ b,\ 0)$ \\
1 & $\xi_2 = -1$ & $(\xi_1, \xi_3)$ & $(a,\ 0,\ b)$ \\
2 & $\xi_1 = +1$ & $(\xi_2, \xi_3)$ & $(P-1,\ a,\ b)$ \\
3 & $\xi_2 = +1$ & $(\xi_1, \xi_3)$ & $(a,\ P-1,\ b)$ \\
4 & $\xi_1 = -1$ & $(\xi_2, \xi_3)$ & $(0,\ a,\ b)$ \\
5 & $\xi_3 = +1$ & $(\xi_1, \xi_2)$ & $(a,\ b,\ P-1)$ \\
\hline
\end{tabular}
\end{center}

\subsubsection{TetGeom}

A nodal distribution of $P(P+1)(P+2)/6$ points, laid out as four vertices,
then the interior of each of the six edges, then the interior of each of the
four faces in \code{faceVerts} order, then the interior of the element:
%
$$
[\, \underbrace{\vec{v}_{0\ldots3}}_{4},\
     \underbrace{e_{0\ldots5}}_{6(P-2)},\
     \underbrace{f_{0\ldots3}}_{4(P-2)(P-3)/2},\
     \underbrace{\text{interior}}_{\text{remainder}} \,].
$$

\paragraph{Trap: edge 2 runs backwards.}  This is the clearest instance of the
modal/nodal divergence of Section
\ref{sec:spatialdomains-orderings-two}.  \code{TetGeom::edgeVerts[2]} is
$\{0,2\}$, but the nodal distribution emits its third edge block running from
vertex $2$ to vertex $0$.  The six blocks are therefore
%
$$
v_0 \! \to \! v_1,\quad v_1 \! \to \! v_2,\quad \mathbf{v_2 \! \to \! v_0},
\quad v_0 \! \to \! v_3,\quad v_1 \! \to \! v_3,\quad v_2 \! \to \! v_3,
$$
%
and that one block must be traversed against the element's local edge
direction.  Unlike the triangle and quadrilateral, nothing is flipped in the
constructor, so \code{m\_eorient[e]} may be used directly for the other five.

\paragraph{Face interiors need aligning.}  A face's nodes are numbered by
whichever element first constructed the \code{TriGeom}, which need not be this
one.  Its interior nodes -- the tail of its curve, after the vertex and edge
blocks -- must therefore be rotated and reflected into this element's view of
the face before being copied.  \code{HOTriangle<T>::Align()} in
\code{SpatialDomains/HOAlignment.h} does this, given the face's own vertex IDs
and the IDs this element expects from \code{faceVerts}.

\subsubsection{PrismGeom}

A nodal distribution of $P(P+1)/2 \cdot P$ points, laid out as six vertices,
then the interior of each of the nine edges, then the interior of each of the
five faces in \code{faceVerts} order -- quadrilateral, triangle,
quadrilateral, triangle, quadrilateral -- and finally the interior of the
element.

\paragraph{Trap: vertices 2 and 3 are transposed.}  The nodal distribution
numbers its vertices in raster order, so nodal slot $2$ sits at
$\vec{\xi} = (-1, 1, -1)$, which is where the standard prism puts vertex $3$.
The permutation from nodal slot to element vertex is therefore
%
$$
\{0,\ 1,\ 3,\ 2,\ 4,\ 5\}.
$$
%
The nine edge blocks, by contrast, do come in \code{edgeVerts} order and each
runs along its edge, so there is no block to reverse as there is for the
tetrahedron.

\paragraph{Mixed faces.}  The triangular faces (1 and 3) are handled as for
the tetrahedron, with \code{HOTriangle::Align()}.  The quadrilateral faces (0,
2 and 4) carry a full tensor grid, so they are first reindexed into the
element's view with \code{ReOrientTracePhysMap} and the interior then taken
with the first face direction fastest.

\subsubsection{PyrGeom}

The pyramid is the remaining gap.  The nodal infrastructure it needs
(\code{StdNodalPyrExp} and an associated point distribution) is newer than
that of the other shapes; once it is available, the layout follows the pattern
above, with the same two warnings as the prism: the nodal distribution numbers
its vertices in raster order, transposing vertices $2$ and $3$, and its edge
blocks $2$ and $3$ are directed $v_3 \! \to \! v_2$ and $v_0 \! \to \! v_3$.

\subsection{Determining an Ordering Empirically}
\label{sec:spatialdomains-orderings-method}

Since neither of the two orderings is authoritative for every shape, the
reliable way to establish the layout for a shape is to read it out of the
point distribution rather than to infer it from a table:

\begin{enumerate}
\item Ask \code{PointsManager()} for the distribution at a modest order --
  order $4$ or so, which is high enough to have a non-trivial number of edge,
  face and interior nodes and low enough to inspect by hand.
\item For each node in turn, classify it by its reference-space position: how
  many of the element's faces it lies on tells you whether it is a vertex, an
  edge node, a face node or an interior node, and which one.
\item Compare the resulting sequence of blocks against the shape's
  \code{edgeVerts} and \code{faceVerts} tables, and record every place the two
  disagree.
\end{enumerate}

The unit test \code{TestMakeOrder.cpp} in \code{library/UnitTests/}
\code{SpatialDomains/} exercises the result.  It is worth noting how it does
so, because the obvious test is vacuous: a single standalone element has
identity face orientations, so its face-reindexing paths are never taken and a
per-shape round-trip test passes whether or not they are correct.  The test
therefore sweeps face orientations explicitly, and each of the reindexing
paths described above has been confirmed to fail the sweep when removed.
